%----------------------------------------------------------------------------------------
% NOTATIE EN FIGUREN VAN DE CURSUS RUIMTEMEETKUNDE
% Staat in de preambule en niet in een hoofdstuk: ook de lijsten met oplossingen gebruiken dit.
%----------------------------------------------------------------------------------------
% Bewaarde grafiek, geplaatst door maak-pdf.py: bestand en pagina in figuren-cache/
\newcommand{\grafiekbestand}[2]{\includegraphics[page=#2,trim=3bp 3bp 3bp 3bp]{figuren-cache/#1}}

\providecommand{\R}{\mathbb{R}}
% Vectoren: \vec{v} met pijl (opmaak), \vc{x}{y}{z} als kolom, \pnt{x}{y}{z} als punt
\newcommand{\vc}[3]{\begin{pmatrix} #1 \\ #2 \\ #3 \end{pmatrix}}
\newcommand{\vcr}[3]{(#1, #2, #3)}
\newcommand{\norm}[1]{\lVert #1 \rVert}
\newcommand{\scal}[2]{#1 \cdot #2}

% Zoals stappen, maar met een bredere kolom voor de berekening: voor lange formules
\newenvironment{stappenbreed}{%
  \par\renewcommand{\arraystretch}{1.25}%
  \noindent\begin{tabular}{@{}>{$\displaystyle}p{9.8cm}<{$}@{\hspace{4mm}}>{\RaggedRight\small\itshape\arraybackslash}p{4.9cm}@{}}%
}{\end{tabular}\par}

% Zoals stappen, met eigen kolombreedtes: \begin{stappenvrij}{6.2cm}{3.8cm}. Voor naast een figuur.
\newenvironment{stappenvrij}[2]{%
  \par\renewcommand{\arraystretch}{1.25}%
  \noindent\begin{tabular}{@{}>{$\displaystyle}p{#1}<{$}@{\hspace{3mm}}>{\RaggedRight\small\itshape\arraybackslash}p{#2}@{}}%
}{\end{tabular}\par}

% Tekst of berekening links, figuur rechts: \naast[breedte rechts]{links}{rechts}
% De breedte rechts komt uit de figuur zelf; het optionele argument wordt niet gebruikt.
\newsavebox{\naastbox}
\newcommand{\naast}[3][]{%
  \par\noindent\sbox{\naastbox}{#3}%
  \begin{minipage}[t]{\dimexpr\linewidth-\wd\naastbox-3mm\relax}\vspace{0pt}#2\end{minipage}\hfill
  \begin{minipage}[t]{\wd\naastbox}\vspace{0pt}\usebox{\naastbox}\end{minipage}\par\addvspace{1mm}}
% Een rij figuren, gecentreerd
\newcommand{\figuren}[1]{\par\addvspace{1mm}{\centering #1\par}\addvspace{1.5mm}}
% Opgave met een grafiek rechts: \begin{metgrafiek}[breedte]{tikz} opgave en deelvragen \end{metgrafiek}
% De solution komt na \end{metgrafiek}, zoals bij metprent.
\newsavebox{\grafiekbox}
\newenvironment{metgrafiek}[2][]{%
  \sbox{\grafiekbox}{#2}%
  \begin{minipage}[t]{\dimexpr\linewidth-\wd\grafiekbox-3mm\relax}%
}{%
  \end{minipage}\hfill
  \begin{minipage}[t]{\wd\grafiekbox}\vspace{-2mm}\usebox{\grafiekbox}\end{minipage}\par}

% Ruimtelijk assenstelsel in schuine projectie: gebruik [ruimte] als optie van tikzpicture,
% en \assendrie{lengte}. De x-as wijst schuin naar voren, y naar rechts, z naar boven.
\tikzset{ruimte/.style={x={(-0.45cm,-0.38cm)},y={(1cm,0cm)},z={(0cm,1cm)}}}
\newcommand{\assendrie}[1]{%
  \draw[->,lijngrijs] (0,0,0) -- (#1,0,0) node[below left,text=black] {$x$};
  \draw[->,lijngrijs] (0,0,0) -- (0,#1,0) node[right,text=black] {$y$};
  \draw[->,lijngrijs] (0,0,0) -- (0,0,#1) node[above,text=black] {$z$};}
